% TOP_PROBLEM: {"id": 9700014, "title": "Second-longest Eulerian excursion on the two-dimensional torus", "category_id": 19, "category": {"id": 19, "name": "probability", "display_name": "Probability", "description": "Problems involving probability theory, stochastic processes, and random structures.", "slug": "probability", "order_index": 19, "created_at": "2026-07-31T15:26:25.671Z"}, "status": "open", "problem_number": "AMR-096-0014", "set_id": 13}
% FIRST_POSTED: 2026-10-01
% ENTRY_KIND: statement_only
% UnsolvedMath ID 9700014; source catalogue code AMR-096-0014
% !TeX program = lualatex
% Definition and sources only; this file is not an LLM solution attempt.
\documentclass[11pt,a4paper]{article}
\usepackage[margin=25mm]{geometry}
\usepackage{fontspec}
\setmainfont{lmroman10-regular.otf}[BoldFont=lmroman10-bold.otf,
 ItalicFont=lmroman10-italic.otf,BoldItalicFont=lmroman10-bolditalic.otf]
\usepackage{amsmath,amssymb,mathtools}
\usepackage{unicode-math}
\setmathfont{latinmodern-math.otf}
\AtBeginDocument{\let\setminus\smallsetminus}
\usepackage{xurl,hyperref}
\hypersetup{colorlinks=true,urlcolor=blue}
\setcounter{secnumdepth}{0}
\setlength{\parindent}{0pt}
\setlength{\parskip}{5pt}
\setlength{\emergencystretch}{3em}
\begin{document}
\section*{Second-longest Eulerian excursion on the two-dimensional torus}\label{9700014}
{\small UnsolvedMath ID 9700014 · AMR-096-0014 · Probability · AMR Open Problem Lists}
\subsection{Definitions and mathematical statement}
For $N\ge3$, take the square torus $G_N=(\mathbb Z/N\mathbb Z)^2$ with nearest-neighbor edges, and replace each edge by one arc in each direction. A directed Eulerian circuit traverses every arc once and returns to its initial vertex. Sample uniformly from these circuits, rooting at the origin; the statistic below is unchanged by cyclic rerooting at another visit to that origin.

The circuit has four excursions from the origin, each running from one departure until the next return. Let their step counts, ordered decreasingly, be
\[
L_1^{(N)}\ge L_2^{(N)}\ge L_3^{(N)}\ge L_4^{(N)}.
\]
Their sum is $4N^2$. Does the random variable $\log L_2^{(N)}/\log N$ converge in distribution as $N\to\infty$ to a random variable $\xi$ whose support is exactly $[0,2]$?

Having that support means $\mathbb P(\xi\in[0,2])=1$ and every nonempty relatively open interval in $[0,2]$ has positive probability. It does not assert that $\xi$ is uniform or has a density. Convergence in distribution may equivalently be tested against every bounded continuous real function.

The second-longest excursion is specified deliberately; the longest must account for a macroscopic amount of the total traversal. The logarithmic scaling asks whether the second-longest length occupies a full range of polynomial scales in $N$, rather than concentrating at a single exponent or at finitely many exponents.
\subsection{Short English statement}
Does the second-longest Eulerian excursion on a square torus have a limiting random growth exponent with full support from zero to two?
\subsection{Sources}
\begin{itemize}
\item[S1] ULAM.AI and the UnsolvedMath Contributors, supplied problems.json, record 9700014 (AMR-096-0014), AMR Open Problem Lists. Catalogue identity and source transcription; CC BY 4.0.
\url{https://huggingface.co/datasets/ulamai/UnsolvedMath}.
\item[S2] Aldous - Open problems, Random Eulerian Circuits, equation (4). Original source indexed by AMR-096-0014.
\url{https://www.stat.berkeley.edu/~aldous/Research/OP/index.html}.
\item[S3] D. Aldous and J. Yu, Random Eulerian Circuits (2014), original numbered conjecture and examples.
\url{https://www.stat.berkeley.edu/~aldous/Papers/aldous_eulerian.pdf}.
\end{itemize}
\end{document}
